% Source: https://artofproblemsolving.com/community/c3h1097217_pythagorean_practice
% Stable identity: aops-topic-1097217-q01
In the sketch (not to scale) the right angle is at $A$. By the Pythagorean theorem,

\begin{asy}
usepackage("amsmath");
size(6cm);
defaultpen(fontsize(10pt));
pair A=(0,0),B=(4.5,0),C=(0,1.5);
draw(A--B--C--cycle,linewidth(0.9));
draw((0.22,0)--(0.22,0.22)--(0,0.22));
label("$A$",A,SW);label("$B$",B,SE);label("$C$",C,NW);
label("$x$",(A+B)/2,S);
label("$9\sqrt{33}$",(A+C)/2,W);
label("$x+1$",(B+C)/2,NE);
label("\mbox{}",(min(currentpicture,user=true)+max(currentpicture,user=true))/2,black); // keeps the text after the figure black when LaTeX typesets the labels inline
\end{asy}

\[
(x+1)^2-x^2=\left(9\sqrt{33}\right)^2 .
\]
The left side is a difference of squares, $(x+1)^2-x^2=2x+1$, and the right side is $81\cdot33=2673$. So $2x+1=2673$, $x=1336$, and the hypotenuse is $x+1=\boxed{1337}$.
